\documentclass{article}
\usepackage[utf8]{inputenc}
\usepackage[russian]{babel}
\usepackage{amsmath}
\usepackage{amssymb}
\usepackage{graphicx}
\usepackage{hyperref}

\title{Пример документа LaTeX}
\author{Автор}
\date{\today}

\begin{document}

\maketitle

\begin{abstract}
Краткое описание содержания документа.
\end{abstract}

\section{Введение}
Это введение к нашему документу. Здесь мы описываем цель и задачи работы. Введение является важной частью любого научного или технического документа, так как оно задаёт тон всему последующему изложению. Хорошее введение должно четко сформулировать проблему, обосновать актуальность исследования и кратко описать структуру документа.

\section{Заключение}
В заключении подводятся итоги проделанной работы. Здесь формулируются основные выводы, полученные в ходе исследования или проекта. Также в заключении можно обсудить возможные направления дальнейшей работы и применения полученных результатов. Заключение должно быть логичным завершением всего документа.

\end{document}